\begin{exercise}{5.5.2}
What is wrong with the following suggestion?
\\
`` The rule 
\[(\Downarrow_{\text{letrec}})
\frac
{M_2[(\text{\textbf{letrec }} x = M_1 \text{\textbf{ in }} M_1)/x] \Downarrow  V, }
{\text{\textbf{letrec }} x = M_1 \text{ in } M_2 \Downarrow  V}
\]
can be simplified to 
\[\frac
{M_2[(\text{\textbf{letrec }} x = M_1 \text{\textbf{ in }} x)/x]\Downarrow V}
{\text{\textbf{letrec}} x = M_1 \text{ in } M_2 \Downarrow  V}
\]
because in the body of the \textbf{letrec}-expression, $x$ id defines to be $M_1$ so we can use \textbf{letrec} $x=M_1$ \textbf{in} $x$ instead of \textbf{letrec} $x = M_1$ \textbf{in} $M_1$.''
\\
Hint: consider \textbf{letrec} $x=0$ \textbf{in} $x$.
\\
\end{exercise}

\begin{solution}
The fundamental error in the proposed simplification of the rule ($\Downarrow_{letrec}$) lies in the fact that it changes the nature of the fixed point of recursion, transforming computations that should terminate into an infinite divergence (non-termination).
In short, while the original rule allows one to “unroll” the recursive definition in order to access the body of the function or expression ($M_1$), the simplified rule replaces the variable with a term that refers back to itself circularly, without ever making progress toward evaluating $M_1$.

\begin{enumerate}
    \item \textbf{Detailed Analysis with the Example: $\textbf{letrec } x = 0 \textbf{ in } 0$} 
    \\
    To show why the suggestion is incorrect, we apply both rules to the expression proposed: $\textbf{letrec } x = 0 \textbf{ in } 0$. \\
    In this case, we have:
    \begin{itemize}
        \item $M_1=0$ (by definition)
        \item $M_2=x$ (by the body of the expression)
    \end{itemize}
    \begin{enumerate}
        \item \textbf{Evaluation Using the Original Rule}\\
        \[(\Downarrow_{\text{letrec}})
\frac
{M_2[(\text{\textbf{letrec }} x = M_1 \text{\textbf{ in }} M_1)/x] \Downarrow  V, }
{\text{\textbf{letrec }} x = M_1 \text{ in } M_2 \Downarrow  V}
\]
\begin{enumerate}
    \item We must evaluate the body $M_2$ (which is $x$), replacing every occurrence of $x$ with the term:
\[\textbf{letrec } x = 0 \textbf{ in } 0\]
    \item The premise expression becomes: 
    \[x[(\textbf{letrec } x = 0 \textbf{ in } 0)/x]\]
    which is
    \[\textbf{letrec } x = 0 \textbf{ in } 0\]
    \item To evaluate $\textbf{letrec } x = 0 \textbf{ in } 0$, we apply the rule again. The new body is $0$. We substitute $x$ in the body, but since $x$ does not occur in $0$, the substitution has no effect.
    \item We then evaluate the final premise: \[0 \Downarrow 0\] which holds by the value axiom.
    \item \textbf{Result}: the expression terminates correctly and returns a value.
\end{enumerate}
        \item \textbf{Evaluation Using the Simplified (Incorrect) Rule}
\[\frac
{M_2[(\text{\textbf{letrec }} x = M_1 \text{\textbf{ in }} x)/x]\Downarrow V}
{\text{\textbf{letrec}} x = M_1 \text{ in } M_2 \Downarrow  V}
\]
\begin{enumerate}
    \item We must evaluate the body $M_2$ (i.e. $x$), replacing every occurrence with:
\[\textbf{letrec } x = 0 \textbf{ in } x\]
\item The premise expression becomes: 
    \[x[(\textbf{letrec } x = 0 \textbf{ in } 0)/x]\]
    which is
    \[\textbf{letrec } x = 0 \textbf{ in } 0\]
    \item To evaluate this premise, we apply the same rule again. 
    \item This forces us to evaluate once more 
    \[\textbf{letrec } x = 0 \textbf{ in } 0\]
    as the premise of the premise.
    \item This generates an infinite chain of identical derivations that never reaches the constant $0$.
    \item \textbf{Result}: the expression diverges and never returns a value.
\end{enumerate}
    \end{enumerate}
	
\item \textbf{Why the Simplification Fails}
\\
The failure comes from a error on how local recursive definitions work:
\begin{itemize}
    \item \textbf{Syntactic Identity vs. Evaluation}:
the term ($\textbf{letrec } x = M_1 \text{\textbf{ in }} x$) is syntactically equivalent to the entire expression we are trying to evaluate.
Using it as a replacement for $x$ means telling the semantic machine: “To know what $x$ is, look at the entire program again.”
This creates an immediate loop.
\item \textbf{Loss of Access to the Definition}: In a $\textbf{letrec }$, $x$ is the name of the recursion, but its value is derived from $M_1$.
Original rule replaces $x$ with a term that points to $M_1$:
\[\textbf{letrec } x = M_1 \text{\textbf{ in }} M_1\]
ensuring that evaluation eventually moves from the variable to the content of the definition.
\item \textbf{Lack of Progress}: A valid inference rule must allow computation to progress toward a normal form or value.
The proposed simplification violates this principle because the premise becomes identical to the conclusion, preventing any reduction.
\end{itemize}

\item \textbf{Conclusion}
\\
The suggestion is wrong because it confuses the \textit{reference} of the recursion ($x$) with the body of the recursion ($M_1$).
By simplifying the rule as proposed, we obtain a language in which no recursive definition can ever be resolved, because every attempt to evaluate the recursive variable results in an infinite self-loop. This makes the $\textbf{letrec } $ construct completely useless for computation.


\end{enumerate}

\end{solution}